\begin{abstract}
eBPF lets developers run custom programs inside the Linux kernel, where a static verifier must prove each one safe before it loads. When the verifier rejects a program, it returns a log of bytecode instructions and register states, and the error flags the operation it could not prove safe. While the log records every step the verifier took, it is usually hard for a developer to tell why the program was rejected or what to change. To identify the gap, we build \empiricalset, a corpus of 235 rejections reproduced from real projects, and conduct an empirical study. We find that the log is hard to use because it does not show two things a repair needs: the root cause and the responsible layer. To address these limitations, we present \sys, a userspace tool that recovers both from the log alone in a single diagnostic. We feed the \sys diagnostic to a downstream repair model and test whether the model repairs the program. Across three models on the 75-case \bench repair suite, \sys raises one-shot repair success by 10.7--21.4 percentage points, from 0.0--37.3\% with the raw log to 10.7--50.7\% with the \sys diagnostic.
\note{1. change tool name (also in the figure); 2. consistency check}
\end{abstract}
